% ============================================================
% ［架空の記入例・提出不可］対象データ，生成条件および評価設計はすべて架空である．
% 通常のchapters/03-problem.texには自動で読み込まれない．
% ============================================================
\chapter{問題設定}

\begin{center}
  \UsamiFictionalExampleNotice
\end{center}

\section{対象データと仮定}

患者$i$の造影腹部CT体積を$x_i\in\mathbb{R}^{H\times W\times D}$，対応する多臓器ラベルを$y_i\in\{0,\ldots,K\}^{H\times W\times D}$とする．クラス0は背景，クラス$1,\ldots,K$は対象臓器を表す．同一患者のスライスまたはパッチが学習，検証，評価の複数集合へ入らないよう，患者単位で集合を固定する．拡散モデルと分割器は学習集合だけで最適化し，ハイパーパラメータおよび合成対の採択閾値は検証集合で決定する．

本例は門脈相に相当する単一の造影相を想定し，CT値の正規化後に共通のボクセル間隔へ再標本化する．臓器ラベルは専門家確認済みであるという架空の前提を置くが，実研究では評価者数，アノテーション手順，評価者間一致度および倫理審査情報を明記する必要がある．また，生成画像は学習用のデータ拡張にだけ使用し，診断画像として提示しない．

\section{多臓器分割}

分割器$f_{\theta}$はCT体積$x_i$を入力し，各ボクセルのクラス確率を出力する．臓器$k$の予測領域を$\hat{Y}_{ik}$，正解領域を$Y_{ik}$とすると，患者$i$，臓器$k$のDice係数を

\begin{equation}
  \operatorname{Dice}_{ik}
  =\frac{2\lvert\hat{Y}_{ik}\cap Y_{ik}\rvert}
  {\lvert\hat{Y}_{ik}\rvert+\lvert Y_{ik}\rvert}
  \label{eq:problem-dice}
\end{equation}

と定義する．主要指標は，評価患者ごとに求めた臓器平均Dice係数とHD95とする．大臓器だけで平均値が決まらないよう臓器別の値も併記し，患者単位の分布から不確実性を示す．式\eqref{eq:problem-dice}の分母が0となる対象は，本例の評価集合には含まれないと仮定する．

\section{合成対の生成と採択}

条件付き生成器$G_{\phi}$は，臓器ラベル$y$と雑音$z$から合成CT $\tilde{x}=G_{\phi}(y,z)$を生成する．学習へ追加する単位は画像だけではなく，合成CTと条件ラベルの対$(\tilde{x},y)$である．各対について，臓器体積の学習分布からの逸脱，左右臓器の重心関係，排他的な臓器間の重複率を検査する．これらをまとめた判定関数$h$により，

\begin{equation}
  b=h(\tilde{x},y)\in\{0,1\}
  \label{eq:problem-synthetic-acceptance}
\end{equation}

を得る．$b=1$の合成対だけを分割器の学習へ加え，$b=0$は除外する．判定関数は臨床的正常を保証するものではなく，あらかじめ定めた範囲から明らかに外れる合成対を除く品質管理として位置づける．評価集合の臓器体積や位置関係を閾値決定に使用しない．

\section{評価課題}

主要比較条件は，実画像のみ，幾何拡張，採択処理を伴わない条件付き拡散拡張，式\eqref{eq:problem-synthetic-acceptance}の採択処理と解剖整合性項を用いる提案法の4条件とする．各条件で分割器の構造，更新回数および評価集合を揃え，データ拡張方法だけを変更する．さらに，採択処理と整合性項を個別に除いた条件を設け，両要素の寄与を確認する．

性能差は平均値だけで判断せず，臓器別Dice，HD95，失敗症例および合成対の採択率を確認する．外部施設での再現性，生成モデルからの学習症例の漏洩，臨床判断への影響は本例の評価範囲外であり，別の検証を要する．本章の症例条件，関数，閾値および比較設計はすべて架空である．
